\begin{pastexam}[label={px:ncku-113-ee-14}]{成大電機所 113 學年度 簡答題第 5 題}
（簡答）

(5 pts) To find biconnected components in a connected graph G, we have to construct a depth
first search (DFS) spanning tree. Please give the reason why the root of the tree must be an
articulation point if it has at least two child.
\end{pastexam}

\begin{pastsoln}
這份考卷和成大電通所 113 學年度的「資料結構」是同一份（電機系、電通所、智慧資安學程共用），完整詳解見歷屆試題~\ref{px:ncku-113-cce-14}。

答案：根的兩棵子樹 $T_1$、$T_2$ 之間沒有邊（無向圖的 DFS 沒有交叉邊，非樹邊都連接祖先和子孫），所以 $c_1$ 到 $c_2$ 的路徑都經過根，刪掉根圖就不連通，根是割點。
\end{pastsoln}
